\documentclass{article}
\usepackage[a4paper,margin=2cm]{geometry}
\usepackage{linguexx}
\begin{document}

\section*{linguexx $\to$ .docx — what to check in Word}

This document was produced by \texttt{linguexx2odt --to docx}. Every example
number below is a live \texttt{SEQ} field and every cross-reference a
\texttt{REF} field, both written with their correct value, so the document
reads right whether or not your reader recalculates on open.

\textbf{1. Does it read correctly?} The examples should number (1) to (5)
in order, and the references at the end should agree with them.

\textbf{2. Does Word offer to repair the file?} It validates against the
ECMA-376 transitional schemas, so a prompt here would be a real finding.

\textbf{3. Does editing renumber?} Put the cursor before example (1), paste
a copy of any example above it, then press Ctrl+A and F9. Everything after
it should move up by one, references included.

\textbf{4. Do the columns hold?} In the glossed examples each gloss word
should sit directly under the word above it. In (5) the label should sit in
a column of its own, level with the object line.

\ex.\label{plain} A plain example, with nothing special about it.

\exg.\label{gloss} il mio libro\\
     the my book\\
\glt `my book'

\ex.
\a. *a judged sub-example
\b. a plain one, whose text should start where the judged one's does

\exg. Der ausserordentlich lange Hund bellte gestern sehr laut.\\
      the extraordinarily long dog barked yesterday very loudly\\
\glt `The extraordinarily long dog barked very loudly yesterday.'

\exg.\label{annot} que Pierre est fatigu\'e \exannot{[CP]}\\
     that Pierre is tired\\
\glt `that Pierre is tired'

Cross-references: \ref{plain}, \ref{gloss}, \ref{annot}, and the bare
form \pref{plain}.

\end{document}
